% A.41 is reserved by the root for the support and logarithm recognition family.
\subsection{A prime-supported coordinate and its opposite cusp}

\begin{statement}{A.41}{Proved}{An odd prime-power neighbour is recognized}
Let $\lambda_0>1$ be an integer.  Suppose that the odd integer among
$\lambda_0,\lambda_0-1$ is a power of one odd prime, and that the
$2$-adic valuation of the even integer is odd.  Every noncuspidal
Selmer section in A.14 with labeled integral realization
$E_{\lambda_0}$ is a rational point section.
In particular, this recognizes the parameters $10,11,14,18,19,23,
24,25,26,27,30,31$ as well as further parameters covered by the
stated arithmetic test.
\end{statement}
\begin{proof}
We first identify the degeneration label from the actual integral
representation.  Let $r$ be an odd prime where $E_{\lambda_0}$ is
multiplicative.  After a common finite extension splitting its
multiplicative torus and the local section's quadratic character,
the rational Tate representation has its unique invariant
cyclotomic line.  Nonzero Tate inertia proves uniqueness: a second
such line or a trivial complement would split the nonzero
extension.  The intersection with the integral $T_2$ is saturated,
so its reduction is a nonzero line in $E[2]$.  In Tate
uniformization it is generated by the point represented by $-1$,
which is the nonidentity 2-torsion point on the smooth part of the
nodal special fibre.  The other two nonzero 2-torsion points are
represented by $\pm q_T^{1/2}$ and specialize to the node; this is
the explicit Tate uniformization of \cite[\S19]{LegIovitaMasdeu2009}.

Write the ordered points as $P_1=(0,0)$, $P_2=(1,0)$,
$P_3=(\lambda_0,0)=P_1+P_2$.  If $r\mid\lambda_0$, the colliding
roots are $0,\lambda_0$, so the integral cyclotomic line reduces to
$\langle P_2\rangle$.  If $r\mid\lambda_0-1$, the colliding roots
are $1,\lambda_0$, so it reduces to $\langle P_1\rangle$.
For a local Legendre parameter tending to infinity, the actual
coordinate change
\[
                x=tX,\qquad y=t\sqrt t\,Y
\tag{A.41.1}
\]
over the indicated quadratic extension gives the Legendre model
with parameter $1/t$; its smooth 2-torsion point corresponds to
$P_3$.  Thus the three contact regions $0,1,\infty$ give respectively
$\langle P_2\rangle,\langle P_1\rangle,\langle P_3\rangle$.
The same calculation applies to a cusp's tangential Tate
representation.  If its coefficient-two tangent inertia were zero,
it would already contradict the realized curve's nontrivial
inertia; otherwise its saturated line gives exactly the stated
nonzero reduction.  The labeled comparison in A.16 reduces to the
identity on $E[2]$, so it preserves these three distinct lines.

Consequently, at an odd prime dividing $\lambda_0$, either a point
centre or a cusp centre has contact label 0; at an odd prime dividing
$\lambda_0-1$ it has contact label 1.  In the ordinary point case
this says respectively $v_r(x_r)>0$ or $v_r(1-x_r)>0$, and the
other coordinate has valuation zero.  At all other odd primes,
good reduction and A.37 give an ordinary integral point.  We have
therefore proved, with the comparison maps just given,
\[
 \begin{split}
  \operatorname{supp}_{\mathrm{odd}}(v_r(A))
     &\subseteq\{r:r\mid\lambda_0\},\\
  \operatorname{supp}_{\mathrm{odd}}(v_r(B))
     &\subseteq\{r:r\mid\lambda_0-1\},
 \end{split}
\tag{A.41.2}
\]
for the completed Kummer classes $A,B$ of $t,1-t$.  Moreover no
finite local cusp has label infinity: at an odd prime the line
calculation excludes it, and at 2 the unequal valuation parities
in the following calculation exclude it.

Put $k=v_2(\lambda_0(\lambda_0-1))$, odd by hypothesis.  If
$\lambda_0$ is even, (A.16.3) gives odd valuation parity for $A_2$
and even parity for $B_2$.  At an ordinary local point,
$v_2(x_2)<0$ would make the two valuations equal; hence
$v_2(x_2)>0$ and $v_2(1-x_2)=0$.  At a cusp, label 1 would force
$A_2=1$ and label infinity would force $A_2/B_2=-1$, both
contradicting these parities.  Only label 0 is possible and it gives
$B_2=1$.  Thus $v_2(B)=0$ in all cases.  If $\lambda_0$ is odd,
exchange $A,B$ to obtain $v_2(A)=0$.

Let the odd integer in the statement be $q^n$.  From (A.41.2) and
the calculation at 2, the coordinate belonging to that integer is
exactly
\[
 C=\begin{cases}
       B=-q^\beta,&\lambda_0\text{ even},\\
       A= q^\beta,&\lambda_0\text{ odd},
    \end{cases}
             \qquad\beta\in\widehat\Z.
\tag{A.41.3}
\]
The signs follow from the global mod-two classes in A.16; the sign
and prime valuations inject into the completed multiplicative
group, as checked in A.39.

Suppose the section were not a rational point.  By A.37 it has at
least two cusp labels.  The preceding paragraph excludes infinity,
so labels 0 and 1 both occur.  Choose a cusp at which $C$ is the
regular coordinate: label 0 if $C=B$, or label 1 if $C=A$.
Write its prime as $p$.  The integral-line comparison shows that
$p\ne q$, and regularity gives $C_p=1$ in
$\widehat{\Q_p^\times}$.  Apply the continuous $p$-adic logarithm
on the unit factor.  It kills the sign and gives
\[
                         \beta_p\log_p(q)=0.
\tag{A.41.4}
\]
Here $\log_p(q)\ne0$.  Indeed the logarithm is injective on
$1+p\Z_p$ for odd $p$, and on $1+4\Z_2$ for $p=2$; its unit
kernel is respectively $\mu_{p-1}$ or $\{\pm1\}$.  A positive
rational prime $q\ne p$ cannot be such a root of unity, since the
corresponding equality $q^{p-1}=1$ or $q^2=1$ would already hold
in $\Q$.  Equation (A.41.4) therefore forces $\beta_p=0$.

We now compare coefficient-$p$ inertia at the distinct odd prime
$q$.  If $s_q$ is an ordinary point, (A.41.2) gives an honest
positive integer valuation for $C_q$.  Its $p$-component cannot be
zero, contradicting $\beta_p=0$.  If $s_q$ is a cusp, its label is
the one where $C$ vanishes, by the saturated-line calculation.
The cusp Tate parameter has leading term
\[
                       q_T=C^2/2^8.
\tag{A.41.5}
\]
More precisely the formal expression is
$C^2/2^8$ times a power series with constant term one; pulling back
to the tangential decomposition group makes that last factor one.
This follows by inverting $j(q_T)=q_T^{-1}+744+\cdots$ in the exact
Legendre formula.  At label 0 the splitting character is that of
$-1$, from the two nodal slopes $\pm\sqrt{-1}$; at label 1 it is
trivial, from slopes $\pm1$.  Both characters are unramified at
the odd prime $q$.  Thus the actual coefficient-$p$ extension
class, after that unramified twist, is
\[
                    2\kappa_p(C)-8\kappa_p(2).
\tag{A.41.6}
\]
The restriction of $\kappa_p(C)$ to inertia at $q$ is
$\beta_p$ times the tame $p$-adic character, whereas 2 is a
$q$-adic unit and has unramified Kummer class.  By (A.41.4) the
representation (A.41.6) is therefore unramified at $q$.

But $E_{\lambda_0}$ has multiplicative reduction at this odd
prime, with $v_q(\Delta)=2n>0$ and unit $c_4$.  Its coefficient-$p$
Tate inertia has nonzero unipotent exponent $2n$, so is ramified.
This contradicts the full integral comparison A.16.  The section
is a point section.  The exact residue computations for the
displayed example parameters are retained in the accompanying
code and output; the proof applies to every integer satisfying the
stated arithmetic test.
\end{proof}
